\chapter{Normas y referencias}
\label{cap:normas}
En este capítulo se pretende recoger las normas y referencias seguidas, así como las herramientas
utilizadas en el desarrollo del proyecto.

\section{Métodos}
\label{sec:metodos}
Para el desarrollo del proyecto se ha seguido una metodología ágil, en concreto Scrum. Se escogió
una metodología ágil debido a que es de gran utilidad para proyectos de desarrollo en los que
pueden surgir cambios en los requisitos o problemas técnicos no previstos que afectarían mucho
a una planificación predictiva \parencite{robert2003agile}. Scrum por su parte es un marco de trabajo
ligero que agrupa el trabajo en ciclos cortos llamados sprints, con duración entre una y cuatro semanas,
donde se planifica el trabajo a realizar y se revisa al final del ciclo \parencite{schwaber2020scrum}.

Este nivel reducido de planificación es más sencillo de realizar y permite un mayor grado de
adaptabilidad a los cambios. Para más información sobre la planificación del proyecto, ver el \tercerAnexo.

\section{Herramientas}
\label{sec:herramientas}

\section{Modelos}
\label{sec:modelos}
Para la detección de posturas en tiempo real se ha utilizado una red neuronal convolucional,
BlazePose, un modelo de aprendizaje profundo diseñado para la detección de posturas humanas en
imágenes y vídeos \parencites{bazarevsky2020blazepose}{xu2020ghum}. Este modelo es capaz de
detectar y rastrear 33 puntos clave (ver figura \ref{fig:blazepose}) del cuerpo humano en tiempo
real, lo que lo hace ideal para aplicaciones de seguimiento de posturas y análisis de movimiento.
Además, incluye junto a las posiciones de los puntos clave en la imagen, la confianza de la
detección y la visibilidad de cada punto, así como una representación de los puntos en 3D.

Las razones para escoger BlazePose sobre otros modelos de detección de posturas se detallan
en el capítulo \textit{\nameref{c08:sec:alternativas-tecnologicas}}, en el apartado dedicado al modelo de
detección de posturas (\ref{c08:subsubsec:mediapipe}). De forma resumida, BlazePose ofrecía un
buen equilibrio entre precisión y velocidad de procesamiento, lo que lo hacía adecuado para esta
aplicación.

\diagrama{
    nombre = {Detección de posturas con BlazePose},
    fuente = {Google, Creative Commons},
    archivo = {memoria/fotos/pl.png},
    etiqueta = {fig:blazepose},
}

\section{Prototipos}
\label{sec:prototipos}
% Diagrama de flujo de la interfaz (UX flow diagram)
\diagrama{
    nombre = {Flujo de la interfaz},
    fuente = {Propia, Creado con Wireflow (\url{https://app.wireflow.co/})},
    archivo = {memoria/fotos/wireflow.jpg},
    width = 0.9\textwidth,
    etiqueta = {fig:wireflow},
}

\section{Métricas}
\label{sec:metricas}
Durante el desarrollo del proyecto era necesario evaluar el correcto funcionamiento de la aplicación,
así como la precisión de la detección de posturas. Aunque para el segundo de estos aspectos el
proceso de validación consistió en probar la aplicación múltiples veces durante el desarrollo de
la misma, para el primero se definieron una serie de pruebas unitarias para garantizar el correcto
funcionamiento de la \gls{api}. Las pruebas unitarias se comenzaron a definir en los primeros sprints
del desarrollo, una vez el se creó la estructura de la \gls{api} y se implementaron las primeras
funciones. Estas pruebas, que se ejecutan de manera casi automática, permiten verificar que cada nueva
versión del código no introduce errores en las funcionalidades ya implementadas.

